% ============================================================
% ch07.tex —— 第 7 章 给运算做体检：群的诞生
% ============================================================
\chapter{给运算做体检：群的诞生}

从这一章起，我们升到高空俯瞰。脚下是已经建成的两个世界：普通数的加减乘（小学工地）、余数王国的加减乘（篇一工地）。眯起眼睛，忽略细节，只看轮廓——你会看到一种\textbf{形状}反复出现。这一章就是给各种运算拍 X 光片，找出所有好运算共享的骨架。这门手艺后来被证明威力惊人：它能一次性解决所有"长得像"的问题，还能预言从未见过的现象。

\section{四项体检}

序章的"$\heartsuit$"运算、普通加法、钟表加法、翻纸牌……这些"运算"质量参差不齐。要比较它们，先得有一张统一的检验单。数学家们（从十八世纪的拉格朗日隐约开始，到十九世纪的伽罗瓦、凯莱明确完成）反复摆弄各种运算之后，提炼出四条最要紧的标准：

\begin{dingyi}[群的公理（体检四项）]
设有一个\textbf{集合} $G$（一堆元素）和一种\textbf{运算} $*$（规则：任取两个元素 $a, b$，能确定地算出一个结果 $a*b$）。逐项检查：
\begin{enumerate}
  \item \Keyword{封闭性}：对任何 $a, b\in G$，$a*b$ 还在 $G$ 里（不出国）；
  \item \Keyword{结合律}：$(a*b)*c = a*(b*c)$（先算哪边不影响结果）；
  \item \Keyword{单位元}：存在一个元素 $e$，对任何 $a$ 都有 $a*e = e*a = a$（什么都不做的元素）；
  \item \Keyword{逆元}：每个 $a$ 都有搭档 $a^{-1}$，使 $a*a^{-1} = a^{-1}*a = e$（人人能回家）。
\end{enumerate}
四项全过，这个"集合 + 运算"的组合就称为一个\textbf{群}（group）。
\end{dingyi}

强调一遍方法论：\textbf{群不是被发明出来烦你的新规定，而是老朋友的共同骨架}。这四条标准不是拍脑袋定的——它们恰好是"能做代数"所需的最低配置：封闭保证算得出来，结合律保证长算式不用打括号（想象一个 $(a*b)*c \ne a*(b*c)$ 的世界：每写三个数就要加括号注明顺序，代数直接瘫痪），单位元和逆元保证"解方程"和"移项"有戏。少任何一条，大厦就缺一根柱子。

\section{体检现场（五位受检者）}

\paragraph{受检者一：整数，加法。}
\begin{itemize}
\item 封闭：整数加整数还是整数 \checkmark；
\item 结合：$(1+2)+3 = 1+(2+3) = 6$，一般显然 \checkmark；
\item 单位元：$0$（$a + 0 = a$）\checkmark；
\item 逆元：$a$ 的搭档是 $-a$（$a + (-a) = 0$）\checkmark。
\end{itemize}
\textbf{通过，是群}。这是你人生中认识的第一个群——只是当时没人发证书。

\paragraph{受检者二：整数，乘法。}
\begin{itemize}
\item 封闭 \checkmark；结合 \checkmark；单位元：$1$（$a\times1 = a$）\checkmark；
\item 逆元：$2$ 的乘法搭档得是 $\tfrac12$（$2\times\tfrac12 = 1$）——\textbf{不是整数}。\textbf{不合格}。
\end{itemize}
补救手术（数学的标准动作：调整会员名单，让体检过关）：
\textbf{方案 A}——会员扩招为"所有非零分数"，逆元 $\tfrac1a$ 到位，通过（这就是分数世界的乘法群）；
\textbf{方案 B}——会员缩编为 $\{1, -1\}$：$1\times1=1$，$1\times(-1) = -1$，$(-1)\times(-1)=1$，人人有逆（各自的逆都是自己）——通过。一个只有两个成员的迷你群，别嫌小，第 8 章它会和"纸片翻面"认亲。

\paragraph{受检者三：减法（作为"运算"本身）。}
结合律直接阵亡：$(5-3)-1 = 1$，而 $5-(3-1) = 3$，不等。\textbf{当场淘汰}，连封闭性都不必查（整数减整数倒还封闭，但一条不合格就出局）。紧接着是数学的标准补刀：\textbf{淘汰减法，把它的活儿交给加法}——$a - b$ 改写成 $a + (-b)$，用"加负数"顶替"减法"。从此数学世界只有一种基本运算（加法）和一个一元操作（取负），减法降格为记号。这与第 6 章"用幂和模乘法顶替大数直接计算"是同一种智慧：\textbf{不留不合格的运算，改造它或解雇它}。

\paragraph{受检者四：钟表世界 $Z_{12}$，加法。}
\begin{itemize}
\item 封闭：余数加余数再压扁，还在 $0$--$11$ 里 \checkmark；
\item 结合：继承普通加法的结合律（压扁只是事后取余，不干扰结合顺序）\checkmark；
\item 单位元：$0$（$a + 0 = a$，绕零圈）\checkmark；
\item 逆元：$7$ 的搭档是 $5$（$7+5 = 12\equiv 0$）；一般地 $a$ 的搭档是 $12-a$（$0$ 自己搭档自己）\checkmark。
\end{itemize}
\textbf{通过}。$Z_n$ 配加法，全体是群——这个家族有无穷多个成员（$n$ 随便取）。

\paragraph{受检者五：序章的"$\heartsuit$"。}
封闭 \checkmark；结合（序章习题你证过：两边都等于 $a+b+c+2$）\checkmark；单位元 $-1$ \checkmark；逆元 $-a-2$ \checkmark。\textbf{通过}。你在序章亲手发明的运算，赫然与两千年历史的整数加法平起平坐——而且序章烧脑时刻你已经看穿：把每个数 $a$ 换成 $a+1$ 记账，$\heartsuit$ \emph{就是}加法穿了件马甲。这件"马甲"有个正式名字，下一章颁发。

\section{最出人意料的受检者：对称也构成群}

前面几位受检者多少都与"数"沾边。现在请一位看起来八竿子打不着的：\textbf{几何动作}。

拿一张等边三角形卡片，问：有几种方式把它\textbf{变动之后仍与原样重合}？（专业说法：保持形状占据同一个位置。转动或翻转都行，只要结束后盖在原来的轮廓上，看不出动过。）

系统地数。旋转类：转 $0^\circ$（没转）、$120^\circ$、$240^\circ$——三种（转 $360^\circ$ 回到 $0^\circ$，不另算；转 $100^\circ$ 之类的，三角形对不上原轮廓，不算）。翻转类：沿三条对称轴各翻一次——三种。合计 \textbf{6 种}。

现在关键一步：定义这些\textbf{动作之间的"乘法"}：
\[
A * B \;:=\; \text{先做动作 } B\text{，再做动作 } A .
\]
（复合。顺序约定先右后左，像穿袜子再穿鞋，脱时先鞋后袜。）逐项体检：
\begin{itemize}
\item 封闭：两个"复原动作"连做，结果仍是复原动作（先各自盖回原样，连做当然也盖回原样）\checkmark；
\item 结合：动作串联天然满足"先做哪两步合并"不影响最终效果——你可以自己拿纸片验证三四个例子，或从"动作是对每个顶点的重排，重排的复合天然可合并"来理解 \checkmark；
\item 单位元：\textbf{"什么都不做"}这个动作 \checkmark（终于，"懒惰"本身成了一个正经的数学元素）；
\item 逆元：每个动作都有反动作——转 $120^\circ$ 的逆是转 $240^\circ$；沿轴翻转的逆是\textbf{再沿同一轴翻一次}（翻过去翻回来）\checkmark。
\end{itemize}
\textbf{对称动作构成群！}这是群论历史上最深刻的一次扩编：\textbf{群不只管数，还管一切"变化"}——镜子、魔方、晶体、分子的振动、方程根之间的换位（第 12 章预告），全部纳入同一张体检表。数学从此有了一台跨领域的翻译机：晶体学家说"这个分子有六重对称"，群论学家听到的就是"这里有个 6 元群"，两拨人隔着专业鸿沟谈笑风生。

\begin{faming}
\textbf{土名}：体检四项、共同骨架 $\to$ \textbf{正式名}：\textbf{群}（group），公理体系由伽罗瓦（第 12 章主角）的工作催生、凯莱 1854 年抽象定型。"抽象代数"这个名字直译了它的本质：抽掉具体对象，只留结构——所以整数加法群、钟表加法群、三角形对称群，可以被\textbf{同一套定理}统治。
\end{faming}

\begin{cidian}
本书土名对照：什么都不做的数/元素 $=$ \textbf{单位元}（identity element）；搭档 $=$ \textbf{逆元}（inverse）；不出国 $=$ \textbf{封闭性}（closure）；能合并先算 $=$ \textbf{结合律}（associativity）；合格运算世界 $=$ \textbf{群}（group）。
\end{cidian}

\section{体检报告能推出什么：两条普遍定理}

公理只有四条，看似空疏，但已经足以推出对所有群都成立的定理——这正是"抽象"的利息：证一次，处处用。

\begin{dingli}[单位元唯一]
一个群只有一个单位元。
\end{dingli}
\begin{proof}
设 $e$ 和 $e'$ 都是单位元。考察乘积 $e * e'$：把 $e'$ 看成单位元（$e$ 看成普通元素），得 $e*e' = e$；把 $e$ 看成单位元（$e'$ 是普通元素），得 $e*e' = e'$。同一个乘积，两个身份，于是 $e = e'$。
\end{proof}
这个证明的味道值得反复咀嚼：它\textbf{什么具体计算都没做}，只让两个候选者互相"演一遍对方的角色"，矛盾自现。这是公理推理的原味。

\begin{dingli}[可约去律]
若 $a*x = a*y$，则 $x = y$（左约去）；若 $x*a = y*a$，则 $x = y$（右约去）。
\end{dingli}
\begin{proof}
左约去：在等式两边\textbf{左乘} $a^{-1}$：
\[
a^{-1}*(a*x) = (a^{-1}*a)*x = e*x = x, \qquad
a^{-1}*(a*y) = (a^{-1}*a)*y = e*y = y .
\]
左边相等推右边相等，$x = y$。（中间一步用了结合律——四条公理在联手干活。）
\end{proof}

第二条正是第 6 章费马小定理证明里"约去 $a^i$"的合法性证书：模素数世界的非零元素构成乘法群（第 10 章正式发证），而\textbf{群里的约去永远安全}——逆元保证你能干净地撤销。\textbf{欠条开始兑现了}：篇一欠下的每个"这里先跳过"，都会在篇二找到它的还款账户。

顺带一条福利推论（习题里你将证明）：逆元也唯一，所以 $a^{-1}$ 这个记号才敢不带标注地写。

\begin{dongshou}
\begin{itemize}
  \yisi 证明：群里每个元素的逆元唯一。（模仿单位元唯一性：设 $b, b'$ 都是 $a$ 的逆元，考察 $b*a*b'$，把 $a*b'$ 先结合得 $b*e = b$，把 $b*a$ 先结合得 $e*b' = b'$，于是 $b = b'$。）
  \yier 给集合 $\{1, 2, 3, 4, 5, 6\}$ 配"乘后模 $7$"做四项体检（这是第 5 章的 $Z_7$ 去掉 $0$）。逐项写一行论证；逆元逐个找：$2^{-1} = 4$（$2\times4 = 8\equiv1$），$3^{-1} = 5$（$15\equiv1$），$6^{-1} = 6$（$36\equiv1$），补全 $1$ 和 $4$ 的。
  \yier 构造一个"只差逆元"的例子：非零整数配乘法（差在哪？）；再构造一个"只差单位元"的例子：正整数配加法（封闭、结合都过，单位元得是 $0$——不在会员名单里）。
  \yisi 用可约去律秒杀一道第 5 章的题：在 $Z_7$（去 $0$）里解 $3x \equiv 2$。两边乘 $3^{-1} = 5$：$x = 10 \equiv 3$。验算 $3\times3 = 9 \equiv 2$ \checkmark。体会"群里的除法 $=$ 乘逆元"。
\end{itemize}
\end{dongshou}

\begin{shaonao}
两条铁律，砸向两个候选者，看谁扛得住。

\textbf{候选一：正有理数乘法群}里解 $x*x = 2$。答案 $x = \sqrt2$——\textbf{不在群里}（无理数）。群公理保证加减乘除畅通，却\textbf{不保证开方畅通}。

\textbf{候选二：$Z_7$ 乘法群}里解 $x*x \equiv 3$。逐个试 $1^2\equiv1$，$2^2\equiv4$，$3^2\equiv2$，$4^2\equiv2$，$5^2\equiv4$，$6^2\equiv1$——$x^2$ 的取值只有 $\{1,2,4\}$，\textbf{$3$ 无解}！

结论：同一个方程 $x^2 = c$，在有的群里永远有解，在有的群里看 $c$ 的脸色。群论管不着开方——想管，得往结构里加更多信息。这个"开方何时有解"的问题，正是伽罗瓦撬动"五次方程没有求根公式"的第一根杠杆（第 12 章选读）：\textbf{解方程的历史，就是研究"开方通行证"的历史}。
\end{shaonao}

\begin{chengshi}
群公理的"极简版"还可以再砍：封闭性可由"运算是 $G\times G\to G$ 的映射"这句话吸收；甚至"左单位元 + 左逆元 + 结合"就够推出全部（教科书习题常客）。本书采用四条版，因为它最贴直觉。另外注意我们默认了"群里的元素个数（阶）可以无穷"——整数加法群就是无穷群；有限群的研究工具（第 9 章的计数论证）不适用无穷群，这是第 21、22 章要绕的弯。
\end{chengshi}

骨架已现，但那位"对称受检者"只亮了相没登场。下一章把它养大：编花名册、画全家福（一张乘法表），然后干一件大事——把一张牌桌上的洗牌，和一张三角形卡片，按在\textbf{同一张}表上。
